% 06-parser-framework.tex — Chapter 6: Parser Framework
\chapter{Parser Framework}
\label{chap:06-parser-framework}
\index{parser!framework}

\pnum
The Lain parser is a hand-written, single-pass, recursive-descent parser.  It is
split across four header files under \srcref{src/parser/}, each responsible for
one category of construct.  The shared infrastructure --- the \cname{Parser}
struct, lookahead management, error reporting, and operator precedence --- lives
in \srcref{src/parser/core.h}.

% -------------------------------------------------------------------------
\section{Parser struct}
\label{sec:par-struct}
\index{Parser struct}
% -------------------------------------------------------------------------

\pnum
The \cname{Parser} struct is the sole mutable state during parsing:

\begin{structdoc}{Parser}{src/parser/core.h:6}
\cname{lexer}  & \cname{Lexer *} & Pointer to the lexer (owns the scan position). \\
\cname{token}  & \cname{Token}   & Current lookahead token (one-token buffer). \\
\cname{line}   & \cname{isize}   & Current source line (incremented on \cname{TOKEN\_NEWLINE}). \\
\cname{column} & \cname{isize}   & Current source column (advanced by token length). \\
\end{structdoc}

\pnum
The parser maintains exactly one token of lookahead in \cname{token}.  Before
any parse function is called, the caller must have already consumed the previous
token so that \cname{token} holds the next unconsumed token.  The invariant is:
\textit{at the entry of every parse function, \texttt{parser->token} is the first
unconsumed token relevant to that function.}

% -------------------------------------------------------------------------
\section{Token advancement and normalization}
\label{sec:par-advance}
\index{parser!token advancement}
% -------------------------------------------------------------------------

\pnum
\cname{\_parser\_advance(parser)} (\srcref{src/parser/core.h:37}) is the single
function that advances the lookahead.  It:
\begin{enumerate}
\item Calls \cname{lexer\_next()} in a loop, discarding
  \cname{TOKEN\_LINE\_COMMENT} and \cname{TOKEN\_MULTILINE\_COMMENT} tokens.
\item Updates \cname{parser->line} and \cname{parser->column} based on the
  raw token.
\item Normalizes \cname{TOKEN\_NEWLINE} to \cname{TOKEN\_EOL}, the canonical
  end-of-line token used by the rest of the parser.
\item Rejects \cname{TOKEN\_SEMICOLON} with error \diagcode{E100}: semicolons
  are not allowed in Lain.
\end{enumerate}

\pnum
The normalization step ensures that all statement-termination logic uses
\cname{TOKEN\_EOL} exclusively.  No parser function ever tests for
\cname{TOKEN\_NEWLINE} or \cname{TOKEN\_SEMICOLON}.

% -------------------------------------------------------------------------
\section{Parser macros}
\label{sec:par-macros}
\index{parser!macros}
% -------------------------------------------------------------------------

\pnum
A family of macros in \srcref{src/parser/core.h:22} provides convenient shorthands:

\begin{center}
\begin{tabular}{ll}
\toprule
\textbf{Macro} & \textbf{Expands to} \\
\midrule
\cname{parser\_match(k)}        & \texttt{parser->token.kind == k} \\
\cname{parser\_advance()}       & \cname{\_parser\_advance(parser)} \\
\cname{parser\_expect(k, msg)}  & error if \texttt{!parser\_match(k)} \\
\cname{parser\_error(msg)}      & unconditional error \\
\cname{parser\_is\_eol()}       & \texttt{parser\_match(TOKEN\_EOL)} \\
\cname{parser\_skip\_eol()}     & advance while EOL or comment \\
\cname{parser\_expect\_eol(msg)}& error if not at EOL \\
\bottomrule
\end{tabular}
\end{center}

\pnum
\cname{parser\_expect(k, msg)} is the most commonly used macro.  It asserts that
the current token has kind \texttt{k}; if not, it calls \cname{\_parser\_error},
which prints the line/column and the message and exits.

% -------------------------------------------------------------------------
\section{Error reporting}
\label{sec:par-errors}
\index{parser!error reporting}
% -------------------------------------------------------------------------

\pnum
Parse errors are reported via \cname{\_parser\_error(parser, msg)}, which
writes to \texttt{stderr} in the format:

\begin{ccode}
[E100] Error Ln %li, Col %li: %message%
\end{ccode}

\pnum
All parse errors use the catch-all code \diagcode{E100}.  Semantic error codes
(E001--E082) are distinct and emitted by sema passes.  After printing, the
function calls \texttt{exit(1)}.  There is no error recovery; the first parse
error terminates the compiler.

% -------------------------------------------------------------------------
\section{Operator precedence}
\label{sec:par-precedence}
\index{operator precedence}
% -------------------------------------------------------------------------

\pnum
Binary operator precedence is defined by \cname{get\_precedence(TokenKind)}
(\srcref{src/parser/core.h:89}).  Higher numbers bind more tightly.  The table
aligns with §8.7 of the Lain language specification:

\begin{center}
\begin{tabular}{rll}
\toprule
\textbf{Level} & \textbf{Operators} & \textbf{Associativity} \\
\midrule
10 & \texttt{* / \%}                                    & Left \\
 9 & \texttt{+ -}                                       & Left \\
 8 & \texttt{<{}<} \quad \texttt{>{}>}                  & Left \\
 7 & \texttt{<} \quad \texttt{<=} \quad \texttt{>} \quad \texttt{>=} \quad \lain{in} & Left \\
 6 & \texttt{==} \quad \texttt{!=}                      & Left \\
 5 & \texttt{\&}                                        & Left \\
 4 & \texttt{\^{}}                                      & Left \\
 3 & \texttt{|}                                         & Left \\
 2 & \lain{and}                                         & Left \\
 1 & \lain{or}                                          & Left \\
\bottomrule
\end{tabular}
\end{center}

\pnum
Operators not in the table return \texttt{-1} from \cname{get\_precedence}, which
terminates the Pratt climbing loop in the expression parser
(see \cref{sec:par-expr-pratt}).

\pnum
The \lain{in} operator is at the same precedence level as the comparison
operators (level~7).  This allows \lain{idx in arr and arr[idx] != 0} to
parse as \lain{(idx in arr) and (arr[idx] != 0)}, where the left-hand \lain{in}
guard proves the right-hand index safe.

% -------------------------------------------------------------------------
\section{Numeric literal parsing}
\label{sec:par-numeric}
\index{numeric literals!parsing}
% -------------------------------------------------------------------------

\pnum
\cname{parse\_numeric\_literal(start, length)} (\srcref{src/parser/core.h:179})
converts a lexer token into a \texttt{long long}.  It strips underscore
separators (e.g.\ \lain{1\_000\_000}), then dispatches on the base prefix:

\begin{itemize}
\item \texttt{0b}/\texttt{0B} prefix: calls \cname{strtoll(buf+2, NULL, 2)}.
\item \texttt{0o}/\texttt{0O} prefix: calls \cname{strtoll(buf+2, NULL, 8)}.
\item \texttt{0x}/\texttt{0X} prefix: calls \cname{strtoll(buf, NULL, 16)}.
\item No prefix: calls \cname{strtoll(buf, NULL, 10)}.
\end{itemize}

\pnum
The function uses a fixed 64-byte buffer for the stripped literal.  A literal
longer than 63 characters after stripping underscores will be silently truncated;
this is acceptable given that all integer types in Lain fit within 64 bits.

% -------------------------------------------------------------------------
\section{Path expression parsing}
\label{sec:par-path}
\index{parser!path expressions}
% -------------------------------------------------------------------------

\pnum
\cname{parse\_path\_expr(arena, parser)} (\srcref{src/parser/core.h:147})
parses a dotted identifier path (e.g.\ \lain{std.io.println}) as a sequence of
\cname{EXPR\_MEMBER} nodes.  It starts with an \cname{EXPR\_IDENTIFIER} and
extends it with member accesses for each dot-separated component.  This helper
is shared between \lain{use} statement parsing and module path parsing.

% -------------------------------------------------------------------------
\section{Parser file structure}
\label{sec:par-files}
% -------------------------------------------------------------------------

\pnum
The parser is split into four files included transitively through \srcref{src/parser.h}:

\begin{center}
\begin{tabular}{ll}
\toprule
\textbf{File} & \textbf{Content} \\
\midrule
\srcref{src/parser/core.h}  & Parser struct, advance, macros, precedence \\
\srcref{src/parser/expr.h}  & Expression parser (Pratt climbing) \\
\srcref{src/parser/stmt.h}  & Statement parser \\
\srcref{src/parser/decl.h}  & Declaration parser (\cname{parse\_module}) \\
\srcref{src/parser/type.h}  & Type expression parser \\
\bottomrule
\end{tabular}
\end{center}
